\section{The largest natural density for rational geometric progressions}
Fix a prime power $q$. We work with monic polynomials over $\mathbb F_q$. A forbidden triple has the form $f,rf,r^2f$ with nonconstant $r\in\mathbb F_q(x)$ and all three terms monic polynomials. In factorization coordinates the three vectors of valuations form an arithmetic progression with a nonzero integral step, which may have either sign in each coordinate. Write $\mathbb P_n$ for the uniform measure on the $q^n$ monic polynomials of degree $n$. The limiting proportion in balls exists with value $b$ exactly when $\mathbb P_n(E)\to b$: subtract successive ball counts for the forward implication and sum the geometric weights for the converse. All results also hold for all nonzero polynomials, since a fixed-scalar class gives an upper bound by the same averaging argument and our constructions may include all scalar multiples.

For $T\ge1$ let $S_T$ consist of the monic irreducibles of degree at most $T$, let $\|h\|=q^{\deg h}$, and set
\[
 c_T=\prod_{P\in S_T}(1-\|P\|^{-1}),\qquad
 v_T=\prod_{P\notin S_T}(1-\|P\|^{-2}).
\]
For subsets $A$ of the monic $S_T$-smooth polynomials containing no forbidden triple define
\[
 B_T=c_T\max_A\sum_{h\in A}\|h\|^{-1}.
\]
The maximum exists. Indeed, identify subsets with their indicator vectors in a countable product of two-point sets. Admissibility is closed, and the objective is continuous by convergence of the full Euler product $1/c_T$.

\paragraph{The extremal theorem.}
If $b$ is the supremum of natural densities of admissible sets, then
\[
 v_T B_T\le b\le B_T,\qquad
 0\le B_T-b\le q^{-T}/(q-1).
\]
In particular $b$ is computable, and there exists an admissible set of natural density exactly $b$.

\paragraph{Upper bound.}
For an admissible set $E$ of density $d$, write $f=gh$ with $g$ $S_T$-rough and $h$ $S_T$-smooth. The fiber $A_g=\{h:gh\in E\}$ is admissible. Average exact-layer proportions for $0\le n<N$. Their sum is
\[
 \sum_{\substack{g\ {\rm rough}\\\deg g<N}}\|g\|^{-1}
 \sum_{\substack{h\in A_g\\\deg h<N-\deg g}}\|h\|^{-1}
 \le (B_T/c_T)\sum_{\substack{g\ {\rm rough}\\\deg g<N}}\|g\|^{-1}.
\]
By finite inclusion--exclusion the rough proportion is exactly $c_T$ in every sufficiently high degree. Division by $N$ and passage to the limit give $d\le B_T$. Applying the sum over all scalar classes proves the same bound when their aggregate has density, whether or not each class individually has density.

\paragraph{Lower bound and computability.}
Use a maximizing $A$ uniformly in every rough fiber for which $g$ is squarefree and omit every other fiber; call this set $E_T$. In an allowed progression the three valuations outside $S_T$ would lie in $\{0,1\}$, so they agree coordinatewise. The common fiber excludes the progression. The density of $E_T$ is $e_T=v_T B_T$. To verify this, first impose a degree bound on $h$ and check squarefreeness only at primes up to a second cutoff. Each resulting condition is a finite Boolean combination of divisor events, and $\mathbb P_n(a\mid f)=\|a\|^{-1}$ once $n\ge\deg a$. The error from large smooth factors is bounded by $\sum_{\deg h>L,\, h\ {\rm smooth}}\|h\|^{-1}$ and that from large squares by $\sum_{\deg P>K}\|P\|^{-2}$. Both bounds hold uniformly in $n$ since the probability of divisibility is always at most the reciprocal norm. The product formula follows on letting the cutoffs grow. Since $B_T\le1$,
\[
 1-v_T\le\sum_{d>T}q^d q^{-2d}=q^{-T}/(q-1).
\]
Finite weighted independent-set optimization on $\deg h\le L$ bounds the unrestricted maximum from below; the preceding smooth tail bounds the deficit. For example the tail is at most $\sum_{n>L}(n+1)^{|S_T|}q^{-n}$. These bounds give an algorithm with any prescribed absolute accuracy.

\paragraph{A uniform estimate for the largest irreducible factor.}
Write $D(f)$ for the largest degree of an irreducible divisor ($D(1)=0$), and $F_m(n)=\mathbb P_n(D(f)\le m)$. For $m\ge1$ we have
\[
 |F_m(n)-F_m(n-1)|\le 1/m\quad(n\ge1).
\]
In fact, let $C(z)=(1-z)\sum_{n\ge0}F_m(n)z^n$. The Euler product over monic polynomials gives
\[
 C(z)=\prod_{\deg P>m}(1-(z/q)^{\deg P})
     =\exp\!\left(-\sum_{k>m}a_kz^k/k\right),\quad
 a_k=q^{-k}\!\sum_{\substack{d\mid k\\d>m}} dI_q(d)\in[0,1],
\]
where $I_q(d)$ counts degree-$d$ monic irreducibles and $\sum_{d\mid k}dI_q(d)=q^k$. If $c_n=[z^n]C(z)$, then $c_n=0$ for $1\le n\le m$. For $n>m$ logarithmic differentiation and coefficientwise majorization by the exponential with positive signs give
\[
 n|c_n|\le\sum_{j=0}^{n-1}|c_j|
 \le\exp\!\left(\sum_{k=m+1}^{n}1/k\right)\le n/m .
\]
The coefficient $c_n$ is precisely the first difference in question. Consequently $|F_m(n-k)-F_m(n)|\le k/m$ for $0\le k\le n$.

\paragraph{Pasting by largest prime degree.}
For each $j\ge1$, take the product set $E_j$ constructed at $T=j$, of density $e_j\to b$, and put $\epsilon_j=2^{-j}$. By the two uniform truncations above, there is a cylinder event $C_j$ with $\mathbb P_n(E_j\mathbin{\triangle} C_j)\le\epsilon_j$ for every $n$. Its indicator has the form $\sum_a u_a\mathbf1_{a\mid f}$ for finitely many monic $a$ by inclusion--exclusion (intersections use least common multiples). Write $s_j=\sum_a u_a/\|a\|$; then $|s_j-e_j|\le\epsilon_j$. Choose increasing positive integers $m_j>j$ exceeding every such $\deg a$ for $C_j$ and satisfying $\sum_a |u_a|\deg(a)/\|a\|\le m_j\epsilon_j$.

Accept $f$ by rule $E_j$ when $m_j<D(f)\le m_{j+1}$, and omit $D(f)\le m_1$. This single set avoids every forbidden triple. Indeed, at a prime of the greatest degree occurring anywhere in a putative triple all three valuations are zero or one: wherever present it is a largest prime and the rule there enforces squarefreeness. As the valuations are in arithmetic progression and one equals one, all equal one. Thus all three polynomials have the same largest degree, follow the same rule, and cannot form the triple.

For $m,n\ge m_j$, uniformity follows from exact division by $a$ and the shift estimate:
\[
 \mathbb P_n(C_j, D\le m)=
 \sum_a (u_a/\|a\|)F_m(n-\deg a)
 =s_jF_m(n)+O(\epsilon_j).
\]
The implicit constant is one. Replacing $C_j$ by $E_j$ and $s_j$ by $e_j$ increases it to three. The joint probability on the $j$th bin therefore differs from $e_j\mathbb P_n(m_j<D\le m_{j+1})$ by at most $6\epsilon_j$; for $n<m_j$ the bin is empty. For each fixed $m$, $F_m(n)\to0$, since the finitely many permitted prime factors give at most $(n+1)^{|S_m|}$ possible polynomials. Hence the total mass of any fixed initial segment of bins tends to zero. The sum of the remaining comparison errors is bounded by $6\sum_{j\ge J}\epsilon_j$, and $\sup_{j\ge J}|e_j-b|\to0$. Letting $n$ and then $J$ tend to infinity proves that the exact-layer proportions tend to $b$.

\paragraph{Effective realization and the greedy comparison.}
The realizing set can be chosen decidable. At stage $j$ use a finite optimizer whose omitted smooth tail, after multiplication by $c_j$, is less than $2^{-j}$. Deterministic tie breaking makes each optimizer computable. The cylinder truncations and the endpoints $m_j$ can then be chosen by the displayed effective bounds. This changes $e_j$ by at most $2^{-j}$ and hence preserves its limit. Factoring a given polynomial and finding its bin decides membership.

Let $G$ be the usual greedy set, consisting of polynomials whose valuations have only $0$ and $1$ as ternary digits. Its full density is the product of its one-prime masses. There is a natural-density improvement over $G$ if and only if for some $T$ there is a \emph{finite} admissible set of $S_T$-smooth polynomials whose reciprocal-norm sum strictly exceeds that of the entire local greedy set. Indeed, a strict local gain can be tensored with the greedy rule at every outside prime. In a progression all outside coordinates must then be constant; convergence of the outside product and its density follow as before, since any excluded exponent is at least two. Conversely, if $b>d(G)$, then for all sufficiently large $T$, $B_T$ exceeds the local greedy mass after normalization by $c_T$, because the omitted greedy product tends to one. Truncating a maximizer in degree preserves a strict inequality for all sufficiently large truncation degrees. Thus strict improvement admits a finite certificate, and exhaustive search will find one whenever it exists. The argument leaves the outcome of this comparison open.

\paragraph{Uniform means on long intervals of degrees.}
The same number $b$ is the largest upper Banach mean of the sequence of exact-degree proportions. In detail, for any admissible $E$ and fixed $T$ one has
\[
 \sum_{n=M}^{M+L-1}\mathbb P_n(E)\le LB_T+O_{q,T}(1)
 \qquad(M\ge0,\ L\ge1).
\]
For rough cores having degrees in $[M,M+L)$ the fiber bound and the eventually exact rough proportions give $LB_T+O_{q,T}(1)$. The contribution of earlier rough cores is at most
\[
 \sum_{s\ge1}\sum_{\substack{h\ {\rm smooth}\\\deg h\ge s}}\|h\|^{-1}
 =\sum_{h\ {\rm smooth}} (\deg h)\|h\|^{-1}<\infty ,
\]
by summing over $s=M-\deg g$ and bounding the rough shell proportion by one. Divide by $L$, take the supremum over the starting position $M$, then let $L$ and $T$ tend to infinity. The reverse extremal inequality follows from the constructed set with converging degree proportions.

\paragraph{A finite binary certificate.}
For $q=2$ choose $S_2$ comprising the two linear primes and the unique quadratic prime. On exponent triples $a+b+2c\le12$ the maximum of $\sum 2^{12-a-b-2c}$ is $14853$, attained by the truncated greedy set. The script \texttt{natural\_density\_check.py} exhausts the independent sets by branching on the earliest remaining vector. Including it deletes the third vertex of any progression whose first two in the chosen ordering have now been included. A branch can be pruned whenever its score plus the total available weight cannot improve the incumbent. All scores are integers. Assigning full mass to the untested smooth tail gives
\[
 b_2\le 1-\frac3{16}\left(\sum_{a+b+2c\le12}2^{-a-b-2c}
                   -\frac{14853}{4096}\right)
 =\frac{2797}{4096}=0.682861328125.
\]
The global greedy density is approximately $0.6483610761$ by the Euler product of Asada et al.
